\documentclass[10pt,a4paper]{amsart}
\usepackage{amsmath,amssymb,amsthm,graphicx}
\usepackage{verbatim}
\newenvironment{reference}{\comment}{\endcomment}
%\newenvironment{reference}{}{}
\newenvironment{slogan}{\comment}{\endcomment}
\newenvironment{history}{\comment}{\endcomment}

% For commutative diagrams we use Xy-pic
\usepackage[all]{xy}

% We use 2cell for 2-commutative diagrams.
\xyoption{2cell}
% Theorem environments.
%
\theoremstyle{plain}
\newtheorem{theorem}[subsection]{Theorem}
\newtheorem{proposition}[subsection]{Proposition}
\newtheorem{lemma}[subsection]{Lemma}

\theoremstyle{definition}
\newtheorem{definition}[subsection]{Definition}
\newtheorem{example}[subsection]{Example}
\newtheorem{exercise}[subsection]{Exercise}
\newtheorem{situation}[subsection]{Situation}

\theoremstyle{remark}
\newtheorem{remark}[subsection]{Remark}
\newtheorem{remarks}[subsection]{Remarks}

\numberwithin{equation}{subsection}

% Macros
%
\def\lim{\mathop{\mathrm{lim}}\nolimits}
\def\colim{\mathop{\mathrm{colim}}\nolimits}
\def\Spec{\mathop{\mathrm{Spec}}}
\def\Hom{\mathop{\mathrm{Hom}}\nolimits}
\def\Ext{\mathop{\mathrm{Ext}}\nolimits}
\def\SheafHom{\mathop{\mathcal{H}\!\mathit{om}}\nolimits}
\def\SheafExt{\mathop{\mathcal{E}\!\mathit{xt}}\nolimits}
\def\Sch{\mathit{Sch}}
\def\Mor{\mathop{\mathrm{Mor}}\nolimits}
\def\Ob{\mathop{\mathrm{Ob}}\nolimits}
\def\Sh{\mathop{\mathit{Sh}}\nolimits}
\def\NL{\mathop{N\!L}\nolimits}
\def\CH{\mathop{\mathrm{CH}}\nolimits}
\def\proetale{{pro\text{-}\acute{e}tale}}
\def\etale{{\acute{e}tale}}
\def\QCoh{\mathit{QCoh}}
\def\Ker{\mathop{\mathrm{Ker}}}
\def\Im{\mathop{\mathrm{Im}}}
\def\Coker{\mathop{\mathrm{Coker}}}
\def\Coim{\mathop{\mathrm{Coim}}}

% Boxtimes
%
\DeclareMathSymbol{\boxtimes}{\mathbin}{AMSa}{"02}

%
% Macros for moduli stacks/spaces
%
\def\QCohstack{\mathcal{QC}\!\mathit{oh}}
\def\Cohstack{\mathcal{C}\!\mathit{oh}}
\def\Spacesstack{\mathcal{S}\!\mathit{paces}}
\def\Quotfunctor{\mathrm{Quot}}
\def\Hilbfunctor{\mathrm{Hilb}}
\def\Curvesstack{\mathcal{C}\!\mathit{urves}}
\def\Polarizedstack{\mathcal{P}\!\mathit{olarized}}
\def\Complexesstack{\mathcal{C}\!\mathit{omplexes}}
% \Pic is the operator that assigns to X its picard group, usage \Pic(X)
% \Picardstack_{X/B} denotes the Picard stack of X over B
% \Picardfunctor_{X/B} denotes the Picard functor of X over B
\def\Pic{\mathop{\mathrm{Pic}}\nolimits}
\def\Picardstack{\mathcal{P}\!\mathit{ic}}
\def\Picardfunctor{\mathrm{Pic}}
\def\Deformationcategory{\mathcal{D}\!\mathit{ef}}
\usepackage{cleveref}
\crefname{originallemma}{original lemma}{original lemmas}
\Crefname{originallemma}{Original lemma}{Original lemmas}
\makeatletter
\newlabel{lemma-recover}{{00VD}{source}{}{sourcepointer-1}{}}
\newlabel{lemma-recover@cref}{{[originallemma][0][]00VD}{[1][0][]source}{}{}{}}
\makeatother
\title{The two functors on presheaves attached to a functor of categories are adjoint}
\author{The Stacks Project Authors}
\date{stacks-project:sites.tex:local-snapshot-20261012; GNU FDL 1.2 or later}
\begin{document}
\maketitle
\noindent\textit{Editorial note.} All delivered mathematical text is the author's own, copied from the pinned Stacks Project chapter file sites-cohomology.tex; no statement or proof step was rewritten, reordered or removed.\par\medskip
\noindent\textit{Editorial note.} The proof is the author's complete construction of the two maps and records the functoriality of each; the sentence 'We will leave it to the reader to verify they are each others inverse' concerns the closing routine verification and is preserved verbatim inside the counted span.\par\medskip
\noindent\textit{Editorial note.} Every labelled reference occurring inside the counted spans is carried in assembly.external\_references as a non-load-bearing navigation pointer (same-chapter ones at their definition site, cross-chapter ones at the invocation inside the counted span), so no load-bearing mathematics is replaced by a stub. No bibliography citation occurs inside the counted spans.\par\medskip
\noindent\textit{Editorial note.} The checked chapter list for same-name collisions is the set of stacks chapters staged in this repository: chow, cohomology, divisors, etale-cohomology, flat, more-algebra, more-morphisms, morphisms, spaces-more-morphisms, varieties together with this card's chapter sites-cohomology.\par\medskip
\noindent\textit{Editorial note.} The preamble is the chapter's own retained minimal preamble (sources/stacks-sites-cohomology/retained-preamble.tex), a verbatim excerpt of the native preamble: the Xy-pic setup and the native theorem declarations with the mathematical macros.\par\medskip
\noindent\textit{Allowed standard prerequisites (editorial).}\begin{itemize}
\item The unit relating a presheaf to the composite of the two functors. Same chapter, Lemma lemma-recover. Supplies the map from a presheaf to the composite used in the first construction.
\item The functors attached to a functor of categories. Same chapter, the definitions of the two functors on presheaves and of the auxiliary category of morphisms into a fixed object. Supplies the functors between which the adjunction is asserted.
\end{itemize}\medskip
\ifdefined\hypertarget\hypertarget{sourcepointer-1}{}\fi
\noindent\textit{Source reference.} Original lemma 00VD: sources/stacks-sites/pinned-original-source.tex, source lines 805--805. Navigation pointer invoked at proof line 805; the counted span contains the pointer only and the cited result is not reproduced inside it, so the pointer is registered as a non-load-bearing navigation stub in the established form for Stacks material. No bibliography citation is involved.\par\medskip
\begin{lemma}
\label{lemma-adjoints-u}
The functor $u_p$ is a left adjoint to the functor $u^p$.
In other words the formula
$$
\Mor_{\textit{PSh}(\mathcal{C})}(\mathcal{F}, u^p\mathcal{G})
=
\Mor_{\textit{PSh}(\mathcal{D})}(u_p\mathcal{F}, \mathcal{G})
$$
holds bifunctorially in $\mathcal{F}$ and $\mathcal{G}$.
\end{lemma}

\begin{proof}
Let $\mathcal{G}$ be a presheaf on $\mathcal{D}$ and let
$\mathcal{F}$ be a presheaf on $\mathcal{C}$.
We will show that the displayed formula holds
by constructing maps either way. We will leave
it to the reader to verify they are each others inverse.

\medskip\noindent
Given a map $\alpha : u_p \mathcal{F} \to \mathcal{G}$
we get $u^p\alpha : u^p u_p \mathcal{F} \to u^p \mathcal{G}$.
Lemma \ref{lemma-recover} says that there is a
map $\mathcal{F} \to u^p u_p \mathcal{F}$. The composition
of the two gives the desired map. (The good thing about this construction
is that it is clearly functorial in everything in sight.)

\medskip\noindent
Conversely, given a map $\beta : \mathcal{F} \to u^p\mathcal{G}$
we get a map $u_p\beta : u_p\mathcal{F} \to u_p u^p\mathcal{G}$.
Let $V \in \Ob(\mathcal{D})$.
We claim that the functor $u^p\mathcal{G}_V$ on $\mathcal{I}_V$
has a canonical map to the constant functor with value $\mathcal{G}(V)$.
Namely, for every object $(U, \phi)$ of $\mathcal{I}_V$,
the value of $u^p\mathcal{G}_V$ on this object is $\mathcal{G}(u(U))$
which maps to $\mathcal{G}(V)$ by $\mathcal{G}(\phi) = \phi^*$.
This is a transformation of functors because $\mathcal{G}$ is a functor
itself. This leads to a map $u_p u^p \mathcal{G}(V) \to \mathcal{G}(V)$.
Another trivial verification shows that this is functorial in $V$
leading to a map of presheaves $u_p u^p \mathcal{G} \to \mathcal{G}$.
The composition $u_p\mathcal{F} \to u_p u^p\mathcal{G} \to
\mathcal{G}$ is the desired map.
\end{proof}

\end{document}
